% ============================================================================
% §2.3 Kết quả — 5 bảng + 2 figure key, ~1.8 trang
% ============================================================================

\subsection{Kết quả}
\label{sec:results}

\subsubsection{Độ tin cậy của đo lường và fit các mô hình LSEM}

Độ tin cậy của \(\hat{\theta}\) được tính theo công thức \(\rho = 1 - \overline{\seit^2}/\widehat{\Var}(\hat\theta)\), đạt 0{,}687 / 0{,}746 / 0{,}814 lần lượt cho ba khối (Bảng~\ref{tab:rho}). Khối~10 nằm ngay dưới ngưỡng thông thường (\(\rho \ge 0{,}70\)) --- đây là cảnh báo khi triển khai các ứng dụng ở cấp cá nhân cho khối này. Tất cả sáu mô hình LSEM (LGCM và LCSM cho ba khối) đều đạt độ phù hợp tốt với CFI \(> 0{,}99\), RMSEA \(< 0{,}031\) và SRMR \(< 0{,}017\) (Bảng~\ref{tab:fit}).

\begin{table}[H]
\centering\small
\caption{Độ tin cậy của ước lượng năng lực ($\rho$) theo khối.}
\label{tab:rho}
\begin{tabular}{lrrr}
\toprule
\textbf{Khối} & \textbf{$n_{\mathrm{obs}}$} & \textbf{$\overline{\seit^2}$} & \textbf{$\rho$} \\
\midrule
10 & 32{.}191 & 0{,}261 & 0{,}687 \\
11 & 29{.}959 & 0{,}295 & 0{,}746 \\
12 & 60{.}807 & 0{,}206 & 0{,}814 \\
\bottomrule
\end{tabular}
\end{table}

\begin{table}[H]
\centering\small
\caption{Fit indices của LGCM và LCSM theo khối.}
\label{tab:fit}
\begin{tabular}{llrrrrr}
\toprule
\textbf{Mô hình} & \textbf{Khối} & \textbf{$\chi^2$} & \textbf{df} & \textbf{CFI} & \textbf{RMSEA} & \textbf{SRMR} \\
\midrule
LGCM & 10 & 17{,}08 & 5 & 0{,}996 & 0{,}019 & 0{,}011 \\
LCSM & 10 &  7{,}11 & 2 & 0{,}998 & 0{,}020 & 0{,}008 \\
LGCM & 11 & 38{,}45 & 5 & 0{,}991 & 0{,}031 & 0{,}017 \\
LCSM & 11 &  6{,}81 & 2 & 0{,}999 & 0{,}019 & 0{,}007 \\
LGCM & 12 & 32{,}50 & 5 & 0{,}997 & 0{,}023 & 0{,}011 \\
LCSM & 12 & 15{,}20 & 2 & 0{,}998 & 0{,}025 & 0{,}010 \\
\bottomrule
\end{tabular}
\end{table}

\subsubsection{Quỹ đạo tăng trưởng và cơ chế đuổi kịp $\to$ plateau}

Hệ số tăng trưởng LGCM dương ở cả ba khối: khối~11 tăng nhanh nhất (\(s = 0{,}028\) logit mỗi đợt, khoảng tin cậy không chứa 0), khối~12 gần như đi ngang (\(s = 0{,}006\)). Hệ số coupling \(\beta\) của LCSM (Hình~\ref{fig:lgcmlcsm}, ô trên-phải) cho thấy một mô hình thay đổi theo khối rất rõ rệt:
\begin{itemize}
  \item Khối~10: \(\beta = -0{,}069\); khoảng tin cậy 95\% \([-0{,}119,\ -0{,}022]\) \textbf{không chứa 0} --- cơ chế đuổi kịp có ý nghĩa thống kê.
  \item Khối~11: \(\beta = -0{,}043\); khoảng tin cậy vẫn giao 0 --- xu hướng đuổi kịp đã suy yếu.
  \item Khối~12: \(\beta = +0{,}002\); khoảng tin cậy bao quanh 0 --- trạng thái ổn định, không còn coupling.
\end{itemize}

Diễn dịch thực tiễn: ở khối~10, một học sinh có \(\theta_{t-1} = -1\) logit (yếu hơn trung bình lớp khoảng một độ lệch chuẩn) được kỳ vọng tăng thêm \(|\beta| \cdot 1 \approx 0{,}07\) logit so với học sinh trung bình ở đợt tiếp theo --- tương đương khoảng 1{,}7\% xác suất trả lời đúng một câu có độ khó trung bình. Hiệu ứng này cộng dồn qua nhiều đợt: sau 5 đợt đánh giá, học sinh yếu được dự đoán đuổi kịp xấp xỉ 0{,}3 logit, tức 7\% xác suất đúng câu trung bình so với học sinh khá cùng lớp.

\begin{figure}[H]
\centering
\includegraphics[width=0.75\linewidth]{main_results.png}
\caption{Bốn tổng quan chính: quỹ đạo LGCM, hệ số coupling \(\beta\) của LCSM, ba ước lượng slope đồng dấu, và tốc độ học tức thời CT-DSEM theo \(\theta\) ở khối 12. Khối 10 = xanh lá, khối 11 = cam, khối 12 = tím.}
\label{fig:lgcmlcsm}
\end{figure}

\subsubsection{Mô hình động vượt baseline tĩnh}

So sánh AIC/BIC và Kiểm định Tỷ số Hợp lý (Likelihood Ratio Test, LRT) giữa mô hình tĩnh (chỉ có intercept) và mô hình động (intercept + slope) khẳng định ưu thế áp đảo của mô hình động (Bảng~\ref{tab:lrt}). \(\Delta\)AIC ở mức 82--167 vượt xa ngưỡng thông thường 10 cho bằng chứng quyết định; giá trị \(p\) của LRT nhỏ hơn \(10^{-19}\) ở cả ba khối và đạt mức \(3 \times 10^{-37}\) ở khối~12.

\begin{table}[H]
\centering\small
\caption{Baseline tĩnh vs động (intercept-only vs intercept + slope ở lmer).}
\label{tab:lrt}
\begin{tabular}{lrrrr}
\toprule
\textbf{Khối} & \textbf{$\Delta$AIC} & \textbf{$\Delta$BIC} & \textbf{LRT $\chi^2$ (df=3)} & \textbf{$p$-value} \\
\midrule
10 &  82 &  57 &  87{,}8 & $6{,}5 \times 10^{-19}$ \\
11 &  92 &  67 &  97{,}7 & $4{,}9 \times 10^{-21}$ \\
12 & 167 & 140 & 172{,}9 & $3{,}0 \times 10^{-37}$ \\
\bottomrule
\end{tabular}
\end{table}

\subsubsection{CT-DSEM khối 12: dynamics chậm về equilibrium}

Trên tập con khối~12 có \(T_i \ge 15\) (\(N = 460\) học sinh), các tham số của mô hình Ornstein--Uhlenbeck (Bảng~\ref{tab:ctdsem}) cho thấy \(\varphi = -0{,}0017\) (khoảng tin cậy âm, dynamics quay về điểm cân bằng), điểm cân bằng dài hạn \(\theta_{\mathrm{eq}} = +0{,}65\) logit (tương ứng \(\approx 66\%\) xác suất trả lời đúng câu trung bình) và bán chu kỳ 408 ngày (vượt một năm học). Tại \(\theta = -1\), tốc độ kỳ vọng đạt \(+0{,}0028\) logit/ngày, tương đương khoảng \(+0{,}25\) logit/học kỳ; tại \(\theta = +1\), tốc độ là \(-0{,}0006\) logit/ngày --- học sinh giỏi có hiện tượng \emph{regression to the mean} nhẹ ngay trong khối~12.

\begin{table}[H]
\centering\small
\caption{Tham số CT-DSEM ở cấp tổng thể (khối 12, $N = 460$, $T_i \ge 15$).}
\label{tab:ctdsem}
\begin{tabular}{lrr}
\toprule
\textbf{Tham số} & \textbf{Ước lượng} & \textbf{CI 95\%} \\
\midrule
$\varphi$ (drift)               & $-0{,}0017$ & $[-0{,}0027, -0{,}0010]$ \\
$\mu \varphi$ (cint)            & $+0{,}0011$ & $[+0{,}0006, +0{,}0016]$ \\
$\sigma$ (diffusion)            & $0{,}046$   & $[0{,}041, 0{,}053]$ \\
manifest var $\sigma_\varepsilon^2$ & $0{,}572$ & $[0{,}562, 0{,}582]$ \\
$\theta_{\mathrm{eq}}$           & $+0{,}65$   & --- \\
half-life (days)                 & 408         & --- \\
\bottomrule
\end{tabular}
\end{table}

\subsubsection{Lợi ích đo lường từ bộ lọc Kalman}

Bộ lọc Kalman cục bộ giúp giảm phương sai sai số đo lường EAP từ 32{,}0\% đến 41{,}2\% mà không cần tăng số câu hỏi khảo sát (Bảng~\ref{tab:ks}, Hình~\ref{fig:kalman}). Phương sai trạng thái ước lượng được \(\hat q_g\) phản ánh biến thiên \emph{thực} của năng lực giữa các đợt đánh giá: khối~11 dao động lớn nhất (\(\sqrt{q} \approx 0{,}86\) logit), phù hợp với giai đoạn học sinh khối 11 chủ động học và làm bài nhiều ngoài giờ trên OLM.

\begin{table}[H]
\centering\small
\caption{Hiệu quả làm trơn của Kalman: SE trung bình giảm, $q_g$ pooled.}
\label{tab:ks}
\begin{tabular}{lrrr r}
\toprule
\textbf{Khối} & \textbf{$\overline{\seit}$ (raw)} & \textbf{$\overline{\seit}$ (smooth)} & \textbf{Giảm (\%)} & \textbf{$\hat q_g$} \\
\midrule
10 & 0{,}499 & 0{,}387 & 41{,}2 & 0{,}444 \\
11 & 0{,}527 & 0{,}433 & 34{,}3 & 0{,}734 \\
12 & 0{,}442 & 0{,}369 & 32{,}0 & 0{,}620 \\
\bottomrule
\end{tabular}
\end{table}

\begin{figure}[H]
\centering
\includegraphics[width=0.68\linewidth]{kalman_overlay_g12.png}
\caption{Chín học sinh khối 12 (ẩn danh). Chấm xám = $\hat\theta$ thô; đường liền = $\hat\theta$ sau Kalman.}
\label{fig:kalman}
\end{figure}

\subsubsection{Ứng dụng: cảnh báo sớm và gợi ý câu hỏi}

Trên quỹ đạo đã làm trơn, tổng tỷ lệ học sinh bị đánh dấu tăng từ khối 10 đến khối 12 (14{,}4\% lên 22{,}3\%; Bảng~\ref{tab:warning}). Nhóm thoả mãn đồng thời cả hai tiêu chí (tốc độ tăng trưởng thấp và bước giảm đột ngột lớn) ổn định ở mức 3{,}5--5{,}7\% --- đây chính là tập học sinh ưu tiên cần giáo viên rà soát thủ công.

\begin{table}[H]
\centering\small
\caption{Số học sinh bị cảnh báo sớm theo khối. $n_{\mathrm{both}}$ là giao của hai tiêu chí.}
\label{tab:warning}
\begin{tabular}{lrrr rrr}
\toprule
\textbf{Khối} & \textbf{$N$ eligible} & \textbf{Total flag} & \textbf{\%} & \textbf{Slope-only} & \textbf{Drop-only} & \textbf{Both} \\
\midrule
10 &  6{.}285 &   905 & 14{,}4 & 411 &   276 & 218 \\
11 &  6{.}740 & 1{.}357 & 20{,}1 & 293 &   683 & 381 \\
12 & 10{.}440 & 2{.}325 & 22{,}3 & 557 & 1{.}281 & 487 \\
\bottomrule
\end{tabular}
\end{table}

Đối với bài toán gợi ý câu hỏi (B7b), 100\% các câu được khuyến nghị có \(|b - \hat\theta_{i,T}| < 0{,}5\) (150 hàng minh hoạ: 10 học sinh đại diện cho mỗi khối, 5 câu cho mỗi học sinh, ba khối). Tỷ lệ này chủ yếu phản ánh \emph{mật độ cao của ngân hàng câu hỏi} (400--800 câu mỗi đơn vị logit) chứ không phải đặc thù riêng của thuật toán; đóng góp mới của bài là khả năng cập nhật gợi ý liên tục theo quỹ đạo đã làm trơn thay vì dựa trên một giá trị \(\hat\theta\) tĩnh duy nhất.
